% RESEARCH_DRAFT_SHA256: 44bb78ffc61a92359f16f68099118b294d5adcb0ef8253889d23ddb48bbf2794
% Layout and source-reference wrappers expanded for the web renderer.
\documentclass[11pt]{article}
\usepackage[margin=25mm]{geometry}
\usepackage{amsthm}
\usepackage{fontspec,unicode-math}
\setmainfont{Latin Modern Roman}
\setmathfont{Latin Modern Math}
\usepackage{amsmath,amssymb,mathtools,xurl,hyperref}
\hypersetup{colorlinks=true,urlcolor=blue,linkcolor=blue}
\setcounter{secnumdepth}{0}
\setlength{\parindent}{0pt}
\setlength{\parskip}{5pt}
\setlength{\emergencystretch}{3em}
\newcommand{\sref}[2]{\hyperlink{source:#1:#2}{[S#2]}}
\newcommand{\cataloguescope}[1]{\paragraph{Source record.}#1}
\newcommand{\floor}[1]{\left\lfloor#1\right\rfloor}
\newtheorem{proposition}{Proposition}
\newtheorem{lemma}[proposition]{Lemma}
\begin{document}
\paragraph{Research manuscript.} OpenAI \texttt{gpt-6-astra}, reasoning \texttt{ultra}; 11 October 2026.
\paragraph{Preserved source body.} The complete catalogue mathematical text and sources follow with layout and source links expanded, including source qualifications and its historical preparation notes. New research begins at the section headed Research attempt.

% UnsolvedMath ID 30005825; source catalogue code OWR-14298166-001
\section{Open Problems in Proof and Computational Complexity}\label{30005825}
{\small UnsolvedMath ID 30005825 · Computer Science}
\subsection{Definitions and mathematical statement}

\paragraph{Shared notation.}Standard number systems are $\mathbb N,\mathbb Z,\mathbb Q,\mathbb R,\mathbb C$. Specialized notation and index ranges follow the cited source and the definitions below.


A refutation is a derivation of contradiction from an unsatisfiable constraint system using specified inference rules. Proof size counts the encoded derivation, and degree in algebraic proof systems is the largest polynomial degree appearing.
\paragraph{Record scope and qualifications.}
The catalogue title is retained as metadata. Named certificates, contracts, bounds or frameworks in that title are not asserted as results of the original source extract.
Source-record limitation: The bundle depends on undefined variants of 3-PIGEON/2-PIGEON and a random-oracle easiness criterion, together with several different proof and graph-density questions; no single complete source formulation can be imported without replacing those definitions.
\paragraph{Statement recorded in the supplied source.}
Is there a total NP search class strictly beyond $\mathrm{FP}$ that is easy in the random-oracle model? Can $3$-PIGEON be randomly black-box reduced to $2$-PIGEON? Are Stabbing Planes and Cutting Planes equivalent up to quasipolynomial simulation? Is nonnegativity of binomials of graph densities decidable? Does the two-sorted theory $\mathrm{VPV}$ prove the nondeterministic time-hierarchy theorem? \href{https://doi.org/10.4171/owr/2024/15}{[S2]}
\paragraph{Source status.}
The supplied catalogue labels this source record open; that metadata was checked against the supplied snapshot on 2026-10-11. The label is preserved for provenance and is not an independent verification of present open status or mathematical correctness. The catalogue literature assessment is dated 2026-08-22.
\subsection{Short English statement}
The source asks about open Problems in Proof and Computational Complexity.
\subsection{Sources}
\begin{itemize}
\item[S1] Ulam AI and the UnsolvedMath Contributors, UnsolvedMath: A Curated Collection of Open Mathematics Problems, record 30005825 (OWR-14298166-001). Catalogue curation and metadata: CC BY 4.0. The source extract is retained; definitions, English summaries and qualifications are editorial additions.
\url{https://huggingface.co/datasets/ulamai/UnsolvedMath}.
\item[S2] Proof Complexity and Beyond (2024), Oberwolfach Reports, DOI 10.4171/owr/2024/15, pp. 871--934
\url{https://doi.org/10.4171/owr/2024/15}.
\item[S3] Oberwolfach Report 15/2024
\url{https://oa.tib.eu/renate/bitstreams/b0ef46dc-1bdf-45b7-a4f5-ca0a9e60a7d6/download}.
\item[S4] Beame et al., Stabbing Planes
\url{https://doi.org/10.4230/LIPIcs.ITCS.2018.10}.
\item Source-specific Oberwolfach Reports licences, including CC BY-SA 4.0 where applicable, are retained; see THIRD\_PARTY\_NOTICES.md.
\end{itemize}
\paragraph{Preparation and review.}
The original source and UnsolvedMath contributors retain attribution for the source extract. Codex (GPT-6; exact serving revision not exposed) assisted with definitions, English summaries, scope notes and formatting. No new human mathematical authorship or independently verified solution is asserted. Permanent source ID 30005825 and code OWR-14298166-001 checked against the complete supplied registry and current repository. Distinct upstream records, including repeated questions and broad or damaged extracts, are retained for complete coverage of the source CS category. A repeated record is not a new independent conjecture. Complete mathematical-equivalence and current-literature review has not been performed. Preparation and final review follow contribution prompts C and D, with source gaps disclosed. Structural and rendering checks do not establish a new theorem or open status.
\label{end:30005825}

\section{Research attempt}
\subsection{Literature, source context and scope}
The full original catalogue body and source qualifications are preserved above. The relevant primary workshop passages were read in the cached publisher report. Historical claims and subsequent attributed results are distinguished from the mathematical work below. A restricted example, standard lemma or known theorem is not counted as new progress on the entire original target. Missing assumptions and repeated or rhetorical extracts are not silently turned into different solved problems.

\subsection{A fully decidable pair inside the graph-density question}
The primary open-problems list contains several independent agendas. We pursue only its binomial homomorphism-density question, retaining the other questions above. Let $G$ be a finite simple graph on $N\ge1$ vertices, and $t(F,G)=\operatorname{Hom}(F,G)/N^{|V(F)|}$; homomorphisms need not be injective. Write $P_3$ for the three-vertex path and $2K_2$ for two disjoint edges.
\begin{proposition}
For a real coefficient $c$, the inequality $c\,t(P_3,G)-t(2K_2,G)\ge0$ holds for every such $G$ if and only if $c\ge1$. The reversed inequality $t(2K_2,G)-c\,t(P_3,G)\ge0$ holds for every $G$ if and only if $c\le0$.
\end{proposition}
\begin{proof}
Writing $d_v$ for the degrees, choose the image of the path center and then independently its two neighbors. Thus
\[
t(P_3,G)=\frac{\sum_v d_v^2}{N^3},\qquad
 t(2K_2,G)=\frac{(\sum_v d_v)^2}{N^4}.
\]
Their difference is $N^{-2}$ times the variance of the degrees under a uniformly chosen vertex, so it is nonnegative. This proves the first inequality for $c\ge1$. In any nonempty regular graph these two densities coincide and are positive, disproving it for $c<1$.
For the reversed inequality, $c\le0$ is immediate from nonnegativity. In the star on $N\ge2$ vertices, the two degree sums are $2(N-1)$ and $N(N-1)$, respectively. Hence
\[
\frac{t(2K_2,G)}{t(P_3,G)}=\frac{4(N-1)}{N^2}\longrightarrow0.
\]
For every $c>0$, a sufficiently large star makes the reversed expression negative.
\end{proof}
This gives an explicit decision rule for the indicated one-parameter families (in particular for rational input coefficients). The attempted general route would express every binomial density difference as a nonnegative degree variance. It already lacks an analogue for arbitrary graph pairs: a path of length two is determined by the degree second moment, whereas a triangle is not. For example, the six-cycle and the disjoint union of two triangles are both 2-regular on six vertices, hence have identical degree moments of every order, but their triangle homomorphism counts are respectively zero and twelve. Therefore degree moments alone cannot decide arbitrary pairs by this route.
\subsection{Verification and remaining gaps}
All simple graphs with at most five vertices were enumerated to compare direct homomorphism counts with both degree formulas; stars and the six-vertex pair were checked separately. This is a standard variance argument in a restricted family, not a general decidability theorem. No conclusion is obtained for the listed TFNP, pigeonhole-reduction, cutting-plane/stabbing-plane, or algebraic-circuit-hierarchy questions.

\section{References and source attribution}
Primary workshop report: \url{https://doi.org/10.4171/owr/2024/15}. Relevant passages read: 2024/15 pp.875--876, five open-problem entries including binomial graph densities.

All retrievals were checked on 11 October 2026. Original question attribution remains in the preserved Sources section. External theorems are attributed and not claimed as new discoveries or fully verified merely by citation.


\paragraph{Integration review.} Codex (GPT-6; exact serving revision not exposed) reviewed the written arguments and expanded layout and source-reference wrappers. Repository metadata and the production web renderer were checked using Python 3.13 and JavaScript with MathJax 3.2.2. This model-assisted review is not independent human proof verification; no PDF compilation or formal proof checking is asserted.

\section{Completion Estimate}
\noindent\textbf{COMPLETION: 0\%}
No substantiated new progress toward the entire original target is claimed. The worked results and failed route are delimited above; a known result, illustrative model, or auxiliary obstruction is not counted as original-target completion.

\section{Final status and attribution}
\textbf{Attempt status: unresolved.} This describes this research attempt, not a new ruling on the source catalogue's status or on external literature claims.
The mathematics and this self-audit were generated by OpenAI \texttt{gpt-6-astra}, reasoning \texttt{ultra}, on 11 October 2026. Codex (GPT-6; exact serving revision not exposed) orchestrated the batch, read and reviewed its written mathematical arguments, and contributed formatting/repository review. No human mathematical authorship, independent human proof checking, or formal verification is claimed. The historical preparation paragraphs in the reproduced catalogue body concern that body only.

Contribution prompt A was applied to the complete definition and supplied source record, followed by the adversarial checks of prompt B. Local reproduction files are \texttt{batch025/checks.py}, \texttt{computations.json}, and the per-ID audit JSON. They were not public links at preparation time. Numerical checks are finite and do not prove unrestricted results. No LaTeX engine was found locally; compilation was not performed by this research agent.

\end{document}
